% !TeX program = xelatex
\documentclass[12pt,a4paper]{article}
\usepackage{fontspec}
\setmainfont{Times New Roman}
\usepackage[vietnamese]{babel}
\usepackage{amsmath,amssymb}
\usepackage[margin=2.2cm]{geometry}
\usepackage{enumitem}
\setlist[enumerate]{itemsep=0.45em,topsep=0.3em}
\begin{document}
\begin{center}{\Large\bfseries ĐÁP ÁN VÀ LỜI GIẢI}\par Đề toán vào lớp 10 Hà Nội, năm học 2023--2024\par\smallskip{\itshape Phần Đại số, toán thực tế và đo lường}\end{center}

\section*{Câu I}
\begin{enumerate}
\item Thay $x=9$ vào $A$:
\[
A=\frac{9+2}{\sqrt9}=\frac{11}{3}.
\]
\item Vì $x>0$, $x\ne1$, ta có
\begin{align*}
B&=\frac{2\sqrt{x}-3}{\sqrt{x}-1}+\frac{3-\sqrt{x}}{(\sqrt{x}-1)(\sqrt{x}+1)}\\
&=\frac{(2\sqrt{x}-3)(\sqrt{x}+1)+3-\sqrt{x}}{(\sqrt{x}-1)(\sqrt{x}+1)}\\
&=\frac{2x-2\sqrt{x}}{(\sqrt{x}-1)(\sqrt{x}+1)}
=\frac{2\sqrt{x}(\sqrt{x}-1)}{(\sqrt{x}-1)(\sqrt{x}+1)}
=\frac{2\sqrt{x}}{\sqrt{x}+1}.
\end{align*}
\item Dùng kết quả trên,
\[
A\cdot B=\frac{x+2}{\sqrt{x}}\cdot\frac{2\sqrt{x}}{\sqrt{x}+1}=\frac{2x+4}{\sqrt{x}+1}.
\]
Điều kiện $A\cdot B=4$ tương đương
\[
2x+4=4\sqrt{x}+4\iff 2\sqrt{x}(\sqrt{x}-2)=0.
\]
Vì $x>0$, suy ra $\sqrt{x}=2$, tức $x=4$ (thỏa điều kiện xác định).
\end{enumerate}

\section*{Câu II}
\begin{enumerate}
\item Gọi $x$ là số sản phẩm theo kế hoạch phân xưởng làm trong một ngày, $x>0$. Theo kế hoạch cần $900/x$ ngày; thực tế mỗi ngày làm $x+15$ sản phẩm và cần $900/(x+15)$ ngày. Thực tế hoàn thành sớm ba ngày nên
\[
\frac{900}{x}-\frac{900}{x+15}=3
\iff x^2+15x-4500=0.
\]
Phương trình có $\Delta=15^2+4\cdot4500=18225=135^2$, nên $x=60$ hoặc $x=-75$. Do $x>0$, nhận $x=60$. Vậy theo kế hoạch mỗi ngày làm 60 sản phẩm.
\item Thể tích hình trụ:
\[
V=\pi r^2h\approx3{,}14\cdot30^2\cdot120=339\,120\ \mathrm{cm}^3.
\]
\end{enumerate}

\section*{Câu III}
\begin{enumerate}
\item Điều kiện $x\ne3$. Nhân phương trình thứ nhất với 3 và phương trình thứ hai với 2:
\[
\begin{cases}
\dfrac6{x-3}-9y=3,\\[3pt]
\dfrac6{x-3}+4y=16.
\end{cases}
\]
Trừ hai phương trình được $13y=13$, do đó $y=1$. Thay vào phương trình đầu: $2/(x-3)-3=1$, suy ra $x-3=1/2$, tức $x=7/2$. Vậy nghiệm là $(x,y)=(7/2,1)$, thỏa điều kiện.
\item Hoành độ giao điểm là nghiệm của
\[
x^2-(m+2)x+m=0.
\]
\begin{enumerate}[label=\alph*)]
\item Biệt thức $\Delta=(m+2)^2-4m=m^2+4>0$ với mọi $m$, nên phương trình có hai nghiệm phân biệt. Do đó $(d)$ luôn cắt $(P)$ tại hai điểm phân biệt.
\item Theo Vi-ét, $x_1+x_2=m+2$ và $x_1x_2=m$. Điều kiện đề bài đòi hỏi các mẫu khác 0, đặc biệt $m\ne0$, và $x_1+x_2-2=m\ne0$. Khi đó
\[
\frac1{x_1}+\frac1{x_2}=\frac{x_1+x_2}{x_1x_2}
=\frac{m+2}{m},\qquad
\frac1{x_1+x_2-2}=\frac1m.
\]
Suy ra $(m+2)/m=1/m$, nên $m+2=1$, tức $m=-1$. Giá trị này thỏa các điều kiện mẫu số; vậy $m=-1$.
\end{enumerate}
\end{enumerate}

\section*{Câu V}
Vì $a,b>0$, các mẫu đều dương. Quy đồng, bất đẳng thức cần chứng minh tương đương
\[
\frac{a^2(b^2+a)+b^2(a^2+b)}{(a^2+b)(b^2+a)}\le1.
\]
Lấy mẫu số trừ tử số:
\begin{align*}
(a^2+b)(b^2+a)-\bigl(a^2(b^2+a)+b^2(a^2+b)\bigr)
&=ab-a^2b^2\\
&=ab(1-ab).
\end{align*}
Từ $a+b\le2$ và $(a-b)^2\ge0$ suy ra $ab\le(a+b)^2/4\le1$. Vì $ab>0$, ta có $ab(1-ab)\ge0$. Do đó hiệu giữa mẫu số và tử số không âm, chứng minh được bất đẳng thức.
\[
1-\left(\frac{a^2}{a^2+b}+\frac{b^2}{b^2+a}\right)
=\frac{ab(1-ab)}{(a^2+b)(b^2+a)}.
\]
\end{document}
